%%%%%%%%%%%%
%
% $Autor: Wings $
% $Datum: 2019-03-05 08:03:15Z $
% $Pfad: MainFunction.tex $
% $Version: 4250 $
% !TeX spellcheck = en_GB/de_DE
% !TeX encoding = utf8
% !TeX root = manual 
% !TeX TXS-program:bibliography = txs:///biber
%
%%%%%%%%%%%%

\chapter{Main Function}

\section{Software Architecture}
The localization system is built on a modular architecture with clearly separated concerns: sensor interfacing, feature extraction, algorithm execution, and result analysis.

\section{Core Modules}
\begin{itemize}
    \item \textbf{Vision Module}: Processes camera frames, extracts features (SIFT, ORB, etc.)
    \item \textbf{Odometry Module}: Computes relative motion from encoders and visual features
    \item \textbf{SLAM Engine}: Performs simultaneous localization and map building
    \item \textbf{Trajectory Manager}: Maintains position history and path records
    \item \textbf{Sensor Fusion}: Integrates multiple sensor streams for robust estimation
\end{itemize}
\begin{figure}[H]
	\centering
	\begin{tikzpicture}[
		node distance=1.5cm,
		module/.style={rectangle, draw=black!70, fill=gray!10, thick, text width=3.5cm, text centered, rounded corners, minimum height=1.5cm},
		arrow/.style={thick, ->, >=stealth}
		]
		\node[module, fill=green!10] (sensor) {\textbf{Sensor Fusion}\\Integrates multiple streams};
		
		\node[module, below left=of sensor] (vision) {\textbf{Vision Module}\\Extracts visual features};
		\node[module, below right=of sensor] (odom) {\textbf{Odometry Module}\\Relative motion estimates};
		
		\node[module, below=of vision, fill=blue!10] (slam) {\textbf{SLAM Engine}\\Simultaneous localization};
		\node[module, below=of odom, fill=red!10] (traj) {\textbf{Trajectory Manager}\\Position history};
		
		\draw[arrow] (sensor) -| (vision);
		\draw[arrow] (sensor) -| (odom);
		\draw[arrow] (vision) -- (slam);
		\draw[arrow] (odom) -- (traj);
		\draw[arrow, dashed] (slam) -- (traj);
	\end{tikzpicture}
	\caption{Core software modules and their data flow within the localization system architecture}
	\label{fig:core-modules}
\end{figure}

\section{Algorithm Pipeline}


\begin{enumerate}
    \item Acquire raw sensor data (image, encoder ticks, optional IMU)
    \item Extract and describe visual features
    \item Match features between consecutive frames
    \item Estimate transformation (motion model)
    \item Update global position estimate
    \item Manage map representation and loop closures
\end{enumerate}

\begin{figure}
	\centering
	\begin{tikzpicture}[
		node distance=0.8cm,
		stepbox/.style={rectangle, draw=blue!60, fill=blue!5, thick, text width=6cm, text centered, rounded corners, minimum height=1cm},
		arrow/.style={thick, ->, >=stealth}
		]
		\node[stepbox] (s1) {\textbf{1. Acquire Data}\\(Image, Encoder ticks, IMU)};
		\node[stepbox, below=of s1] (s2) {\textbf{2. Extract Features}\\(SIFT, ORB, etc.)};
		\node[stepbox, below=of s2] (s3) {\textbf{3. Match Features}\\Between consecutive frames};
		\node[stepbox, below=of s3] (s4) {\textbf{4. Estimate Transformation}\\(Motion model)};
		\node[stepbox, below=of s4] (s5) {\textbf{5. Update Global Position}\\Global position estimate};
		\node[stepbox, below=of s5] (s6) {\textbf{6. Map Management}\\Manage map \& loop closures};
		
		\draw[arrow] (s1) -- (s2);
		\draw[arrow] (s2) -- (s3);
		\draw[arrow] (s3) -- (s4);
		\draw[arrow] (s4) -- (s5);
		\draw[arrow] (s5) -- (s6);
	\end{tikzpicture}
	\caption{Six-stage localization algorithm pipeline, from raw sensor acquisition to map and loop-closure management}
	\label{fig:algo-pipeline}
\end{figure}